% GENERATED FILE. Edit canonical lessons, then run npm run generate:books.
% canonical-source-hash: fnv1a64:a4ba68ca855a5329
% canonical-lessons: RU-C133-to-go-up, RU-C133-to-go-down, RU-C133-to-stay, RU-C133-to-rent, RU-C133-to-repair

\chapter{To go up, To go down, To stay, To rent, To repair}
\label{ch:ru-to-go-up-to-go-down-to-stay-to-rent-to-repair}
\hlchaptermodality{133}

\begin{chapteropening}
\textbf{By the end of this chapter:} \emph{I can say to go up, to go down, to stay, to rent and to repair.}

Complete the last lesson of chapter 133: I can say to go up, to go down, to stay, to rent and to repair.
\end{chapteropening}

\section[\texorpdfstring{\ru{подниматься} (podnimátsya)}{подниматься (podnimátsya)}]{\ru{подниматься} (podnimátsya) — to go up}

\label{lesson:RU-C133-to-go-up}



\begin{quote}
Before the new one: say the Russian for to sit down, then the Russian for to fall.
\end{quote}

\subsection*{You'll want to know: \ru{подниматься}}
\textbf{\ru{подниматься}} (\emph{podnimátsya}) — \textquotedblleft{}to go up\textquotedblright{}.

\begin{cousinweb}
\textbf{\ru{подниматься}} is \emph{\ru{под}-}, \textquotedblleft{}up from under\textquotedblright{}, with \emph{-\ru{нимать}}, from the old verb \emph{*jęti}, \textquotedblleft{}to take\textquotedblright{}, and \emph{-\ru{ся}}, \textquotedblleft{}oneself\textquotedblright{}: to lift yourself up.
\end{cousinweb}

\begin{grammarlens}[title={What you've built}]
One more word for this chapter.
\end{grammarlens}

\subsection*{Your turn}
\begin{itemize}
\raggedright
  \item \emph{Say it:} \emph{podnimátsya}
  \item \emph{Say it:} \emph{podnimátsya}, once more
  \item \emph{Say it:} say \emph{podnimátsya}
  \item \emph{Recall:} say the Russian for to sit down, then the Russian for to fall, then say \emph{podnimátsya} again
\end{itemize}

\begin{tcolorbox}[breakable,colback=teal!4,colframe=teal!35!black,title={Before you move on}]
What does \ru{подниматься} mean? (\textquotedblleft{}To go up\textquotedblright{}.) Say it once more.
\end{tcolorbox}
\section[\texorpdfstring{\ru{спускаться} (spuskátsya)}{спускаться (spuskátsya)}]{\ru{спускаться} (spuskátsya) — to go down}

\label{lesson:RU-C133-to-go-down}



\begin{quote}
Before the new one: say the Russian for to fall, then the Russian for to go up.
\end{quote}

\subsection*{You'll want to know: \ru{спускаться}}
\textbf{\ru{спускаться}} (\emph{spuskátsya}) — \textquotedblleft{}to go down\textquotedblright{}.

\begin{cousinweb}
\textbf{\ru{спускаться}} is \emph{\ru{с}-}, \textquotedblleft{}down\textquotedblright{}, with \emph{\ru{пускать}}, \textquotedblleft{}to let go\textquotedblright{}, and \emph{-\ru{ся}}, \textquotedblleft{}oneself\textquotedblright{}: to let yourself down. Vasmer links \emph{\ru{пускать}} with \emph{\ru{пустить}} and \emph{\ru{пустой}}, \textquotedblleft{}empty\textquotedblright{}; its deeper origin is uncertain.
\end{cousinweb}

\begin{grammarlens}[title={What you've built}]
One more word for this chapter.
\end{grammarlens}

\subsection*{Your turn}
\begin{itemize}
\raggedright
  \item \emph{Say it:} \emph{spuskátsya}
  \item \emph{Say it:} \emph{spuskátsya}, once more
  \item \emph{Say it:} say \emph{spuskátsya}
  \item \emph{Recall:} say the Russian for to fall, then the Russian for to go up, then say \emph{spuskátsya} again
\end{itemize}

\begin{tcolorbox}[breakable,colback=teal!4,colframe=teal!35!black,title={Before you move on}]
What does \ru{спускаться} mean? (\textquotedblleft{}To go down\textquotedblright{}.) Say it once more.
\end{tcolorbox}
\section[\texorpdfstring{\ru{оставаться} (ostavátsya)}{оставаться (ostavátsya)}]{\ru{оставаться} (ostavátsya) — to stay}

\label{lesson:RU-C133-to-stay}



\begin{quote}
Before the new one: say the Russian for to go up, then the Russian for to go down.
\end{quote}

\subsection*{You'll want to know: \ru{оставаться}}
\textbf{\ru{оставаться}} (\emph{ostavátsya}) — \textquotedblleft{}to stay\textquotedblright{}.

\begin{cousinweb}
\textbf{\ru{оставаться}} is \emph{\ru{о}-} with \emph{-\ru{ставать}}, from \emph{\ru{стать}}, \textquotedblleft{}to stand\textquotedblright{}, and \emph{-\ru{ся}}: to remain where you stand.
\end{cousinweb}

\begin{grammarlens}[title={What you've built}]
One more word for this chapter.
\end{grammarlens}

\subsection*{Your turn}
\begin{itemize}
\raggedright
  \item \emph{Say it:} \emph{ostavátsya}
  \item \emph{Say it:} \emph{ostavátsya}, once more
  \item \emph{Say it:} say \emph{ostavátsya}
  \item \emph{Recall:} say the Russian for to go up, then the Russian for to go down, then say \emph{ostavátsya} again
\end{itemize}

\begin{tcolorbox}[breakable,colback=teal!4,colframe=teal!35!black,title={Before you move on}]
What does \ru{оставаться} mean? (\textquotedblleft{}To stay\textquotedblright{}.) Say it once more.
\end{tcolorbox}
\section[\texorpdfstring{\ru{снимать} (snimát)}{снимать (snimát)}]{\ru{снимать} (snimát) — to rent; to take off}

\label{lesson:RU-C133-to-rent}



\begin{quote}
Before the new one: say the Russian for to go down, then the Russian for to stay.
\end{quote}

\subsection*{You'll want to know: \ru{снимать}}
\textbf{\ru{снимать}} (\emph{snimát}) — \textquotedblleft{}to rent; to take off\textquotedblright{}.

\begin{cousinweb}
\textbf{\ru{снимать}} is \emph{\ru{с}-}, \textquotedblleft{}off\textquotedblright{}, with \emph{-\ru{нимать}}, from the old verb \emph{*jęti}, \textquotedblleft{}to take\textquotedblright{}: \textquotedblleft{}to take off\textquotedblright{}, and so \textquotedblleft{}to take\textquotedblright{} a flat for rent.
\end{cousinweb}

\begin{grammarlens}[title={What you've built}]
One more word for this chapter.
\end{grammarlens}

\subsection*{Your turn}
\begin{itemize}
\raggedright
  \item \emph{Say it:} \emph{snimát}
  \item \emph{Say it:} \emph{snimát}, once more
  \item \emph{Say it:} say \emph{snimát}
  \item \emph{Recall:} say the Russian for to go down, then the Russian for to stay, then say \emph{snimát} again
\end{itemize}

\begin{tcolorbox}[breakable,colback=teal!4,colframe=teal!35!black,title={Before you move on}]
What does \ru{снимать} mean? (\textquotedblleft{}To rent; to take off\textquotedblright{}.) Say it once more.
\end{tcolorbox}
\section[\texorpdfstring{\ru{чинить} (chinít)}{чинить (chinít)}]{\ru{чинить} (chinít) — to repair}

\label{lesson:RU-C133-to-repair}



\begin{quote}
Before the new one: say the Russian for to stay, then the Russian for to rent; to take off.
\end{quote}

\subsection*{You'll want to know: \ru{чинить}}
\textbf{\ru{чинить}} (\emph{chinít}) — \textquotedblleft{}to repair\textquotedblright{}.

\begin{cousinweb}
\textbf{\ru{чинить}} comes from \emph{\ru{чин}}, \textquotedblleft{}order, rank\textquotedblright{} (Vasmer): to repair is to put back in order. Vasmer relates \emph{\ru{чин}} to Sanskrit \emph{cinóti}, \textquotedblleft{}arranges, piles\textquotedblright{}, and Greek \emph{poiéō}, \textquotedblleft{}I make\textquotedblright{}.
\end{cousinweb}

\begin{grammarlens}[title={What you've built}]
One more word for this chapter.
\end{grammarlens}

\subsection*{Your turn}
\begin{itemize}
\raggedright
  \item \emph{Say it:} \emph{chinít}
  \item \emph{Say it:} \emph{chinít}, once more
  \item \emph{Say it:} say \emph{chinít}
  \item \emph{Recall:} say the Russian for to stay, then the Russian for to rent; to take off, then say \emph{chinít} again
\end{itemize}

\begin{tcolorbox}[breakable,colback=teal!4,colframe=teal!35!black,title={Before you move on}]
What does \ru{чинить} mean? (\textquotedblleft{}To repair\textquotedblright{}.) Say it once more.
\end{tcolorbox}
